% from ega4-2-fr.tex
\oldpage[IV-2]{163}

Remplaçant au besoin $A$ par un anneau de fractions $A_f$ (avec $f\ne0$),
on peut en outre supposer que $A'$ contient un anneau de polynômes
$A[T_1,\ldots,T_n]=A''$, et est une $A''$-algèbre finie (Bourbaki,
\emph{Alg. comm.}, chap. V, \S3, n\textsuperscript{o} 1, cor. 1 du th. 1).
Mais $A''$ est un anneau de Cohen--Macaulay (6.3.6); donc on peut se borner
au cas où de plus $A'$ est une $A$-algèbre finie. Il y a alors $g\ne0$ dans
$A$ tel que $A'_g$ soit un $A_g$-module libre de type fini (Bourbaki,
\emph{Alg. comm.}, chap. II, \S5, n\textsuperscript{o} 1, cor. de la
prop. 2), donc on peut encore supposer de plus que $A'$ soit un $A$-module
libre. Mais alors $A'$ est un $A$-module de Cohen--Macaulay (0, 16.5.1), et
comme $A'$ est un $A$-module de type fini, $A'$ est aussi un $A'$-module de
Cohen--Macaulay (0, 16.5.3), donc un anneau de Cohen--Macaulay.
\item[(ii)] Supposons qu'il existe un $\mathcal O_X$-Module cohérent
$\mathcal F$ tel que $\operatorname{Supp}(\mathcal F)=X$ et que
$\mathcal F$ soit un $\mathcal O_X$-Module de Cohen--Macaulay. Alors $X$
vérifie la condition (CMU): en effet, avec les notations de la démonstration
de (6.11.8), la relation (6.11.8.1) montre que l'on a
$\operatorname{coprof}(\mathcal O_{Y,x})=0$ dans un voisinage (par rapport
à $Y$) du point générique $y$ de $Y$.
\end{enumerate}

On ignore s'il existe des préschémas localement noethériens de dimension
$\geq2$ qui ne vérifient pas (CMU) (si $\dim(X)=1$, il est immédiat que tout
point maximal de $X_{\mathrm{red}}$ admet un voisinage ouvert intègre de
dimension $1$, donc de Cohen--Macaulay).
\end{remarks}

\subsection{\texorpdfstring{Critères de Nagata pour que $\operatorname{Reg}(X)$ soit ouvert}{Critères de Nagata pour que Reg(X) soit ouvert}}
\label{IV.2.s6.12-fr}

\begin{env}[6.12.1]
\label{IV.2.s6.12.1-fr}
Étant donné un préschéma localement noethérien $X$, on appelle lieu
singulier de $X$ et on note $\operatorname{Sing}(X)$ l'ensemble des points
$x\in X$ tels que $X$ ne soit pas régulier au point $x$ (autrement dit,
tels que l'anneau local $\mathcal O_{X,x}$ ne soit pas régulier); le
complémentaire $X-\operatorname{Sing}(X)$, ensemble des $x\in X$ où $X$ est
régulier, se note $\operatorname{Reg}(X)$. Nous nous proposons dans ce
numéro de donner des conditions moyennant lesquelles
$\operatorname{Sing}(X)$ est fermé (i.e. $\operatorname{Reg}(X)$ ouvert).
\end{env}

\begin{proposition}[6.12.2]
\label{IV.2.s6.12.2-fr}
Soit $X$ un préschéma localement noethérien. Les conditions suivantes sont
équivalentes:
\begin{enumerate}
\item[a)] $\operatorname{Reg}(X)$ est ouvert dans $X$.
\item[b)] Pour tout $x\in\operatorname{Reg}(X)$, il existe une partie
ouverte non vide dans $\overline{\{x\}}$ contenue dans
$\operatorname{Reg}(X)$.
\end{enumerate}
En outre, ces conditions sont entraînées par la suivante:
\begin{enumerate}
\item[c)] Pour tout $x\in\operatorname{Reg}(X)$, si on désigne par $Y$ le
sous-préschéma fermé réduit de $X$ ayant $\overline{\{x\}}$ pour espace
sous-jacent, $\operatorname{Reg}(Y)$ est un voisinage de $x$ dans $Y$.
\end{enumerate}
\end{proposition}

\begin{proof}
L'équivalence de a) et b) résulte de ce que $\operatorname{Reg}(X)$ est
stable par générisation (0, 17.3.2) et de (0$_{\rm III}$, 9.2.6). Pour
prouver que c) entraîne b), on peut se borner au cas où
$X=\operatorname{Spec}(A)$ est un voisinage ouvert affine de $x$; si
$(t_i)_{1\leq i\leq n}$ est un système régulier de paramètres (0, 17.1.6)
de l'anneau local régulier $A_x$, on peut supposer (en remplaçant au besoin
$X$ par un voisinage ouvert de $x$) que $t_i=(s_i)_x$ pour tout $i$, où
$s_i\in A$ et où la famille $(s_i)_{1\leq i\leq n}$ est $A$-régulière
(0, 15.2.4). On a alors $Y=\operatorname{Spec}(A/\mathfrak p)$, où
$\mathfrak p=\mathfrak j_x$; comme les $t_i$ engendrent l'idéal maximal
$\mathfrak m=\mathfrak pA_x$ de $A_x$, on peut encore supposer (en
remplaçant $X$ par un voisinage plus petit de $x$) que les $s_i$ engendrent
$\mathfrak p$. Pour tout $y\in Y$, les $(s_i)_y$ engendrent donc
$\mathfrak p_y$; comme ils forment une suite $A_y$-régulière, ils font
partie d'un système de paramètres de $A_y$ (0, 16.4.1); donc on déduit de
(0, 17.1.7) que si $(A/\mathfrak p)_y=A_y/\mathfrak p_y$ est régulier, il
en est de même de $A_y$, d'où la conclusion.
\end{proof}

\begin{corollary}[6.12.3]
\label{IV.2.s6.12.3-fr}
Soit $X$ un préschéma localement noethérien. Les conditions suivantes sont
équivalentes:
\begin{enumerate}
\item[a)] Pour tout sous-préschéma $Y$ de $X$,
$\operatorname{Reg}(Y)$ est ouvert dans $Y$.

